% hep_COT — 讲清楚"这项工作做什么 + 为什么有意义"的流程图
% Compile: xelatex pipeline_diagram.tex
\documentclass[11pt]{article}
\usepackage[a3paper, landscape, margin=1.2cm]{geometry}
\usepackage[UTF8]{ctex}
\usepackage{xcolor}
\usepackage{tikz}
\usepackage{amsmath}
\usepackage{amssymb}
\usepackage{array}
\usepackage{adjustbox}
\usetikzlibrary{
  positioning, arrows.meta, shapes.geometric, shapes.misc,
  calc, fit, backgrounds, decorations.pathreplacing, decorations.markings
}

% --- 配色：以"意义"为主，弱化代码感 ---
\definecolor{cPaper}{RGB}{70, 90, 120}          % 输入：论文
\definecolor{cPhase}{RGB}{60, 130, 180}         % 六维分析
\definecolor{cProduct}{RGB}{60, 140, 100}       % 产出
\definecolor{cTeacher}{RGB}{60, 130, 180}       % teacher
\definecolor{cNative}{RGB}{60, 160, 110}        % native
\definecolor{cAligned}{RGB}{220, 130, 50}       % aligned
\definecolor{cSignal}{RGB}{140, 90, 170}        % 反馈信号（规则化 judge）
\definecolor{cInsight}{RGB}{180, 60, 90}        % 意义/价值
\definecolor{cBG}{RGB}{247, 247, 250}

\tikzset{
  base/.style={
    rectangle, rounded corners=3pt, draw=black!40, line width=0.6pt,
    align=center, font=\small, inner sep=7pt, minimum height=9mm
  },
  paper/.style   ={base, fill=cPaper!12,   draw=cPaper!70,   text=cPaper!85!black},
  phase/.style   ={base, fill=cPhase!10,   draw=cPhase!60,   text=cPhase!85!black,
                   text width=42mm, minimum height=20mm},
  product/.style ={base, fill=cProduct!13, draw=cProduct!75, text=cProduct!85!black},
  teacher/.style ={base, fill=cTeacher!13, draw=cTeacher!70, text=cTeacher!85!black},
  native/.style  ={base, fill=cNative!15,  draw=cNative!75,  text=cNative!85!black},
  aligned/.style ={base, fill=cAligned!20, draw=cAligned!80, text=cAligned!85!black},
  signal/.style  ={base, fill=cSignal!14,  draw=cSignal!75,  text=cSignal!85!black},
  insight/.style ={base, fill=cInsight!10, draw=cInsight!75, text=cInsight!85!black,
                   font=\small},
  note/.style={font=\footnotesize\itshape, text=black!60, align=left},
  arrow/.style={-{Stealth[length=3mm, width=2.4mm]}, line width=0.7pt, draw=black!70},
  fatarrow/.style={arrow, line width=1.2pt, draw=cInsight!80},
  darrow/.style={arrow, dashed, draw=black!55},
  loop/.style={arrow, draw=cAligned!85, line width=0.85pt},
  signalarrow/.style={arrow, draw=cSignal!80, line width=0.85pt},
  title/.style={font=\LARGE\bfseries},
  subtitle/.style={font=\normalsize, text=black!70},
  banner/.style={rectangle, rounded corners=6pt, fill=cBG, draw=black!20,
                 line width=0.4pt, inner sep=10pt, align=left},
}

\begin{document}
\pagestyle{empty}

% ===================================================================
% PAGE 1 — 这项工作给一篇 HEP 论文"产出"了什么
% ===================================================================
\begin{center}
{\LARGE\bfseries hep\_COT：让大语言模型"读懂"一篇 HEP 实验论文}\\[2pt]
{\small 把一篇实验论文压成一份\textbf{结构化、可机读、可审计}的六维深度报告}
\end{center}

\vspace{2mm}

\begin{adjustbox}{max totalsize={\textwidth}{0.86\textheight}, center}
\begin{tikzpicture}[node distance=8mm and 14mm]

% ---------- 左：输入 ----------
\node[paper, text width=56mm] (paper)
  {\textbf{输入：一篇 HEP 实验论文}\\[2pt]
   \footnotesize arXiv tex 源码\\
   \footnotesize $+$ 图表 (pdf/png + caption)\\
   \footnotesize $+$ HEPData 数据表\\
   \footnotesize $+$ INSPIRE 引用图\\[4pt]
   \footnotesize\itshape 例：CMS VBS same-sign WWH \\ \itshape\footnotesize (arXiv 2206.08956)};

% ---------- 中：六维分析（叙事化，不提 schema key） ----------
\node[subtitle, right=16mm of paper, yshift=22mm] (phasetitle)
  {\textbf{一次完整分析 = 六个物理问题，逐一作答}};

% 用 2x3 网格摆六个 phase
\node[phase, below=4mm of phasetitle, xshift=-48mm] (p1)
  {\textbf{1. 测了什么}\\[2pt]
   \footnotesize 测量类型、终态、能量与亮度、\\
   \footnotesize 头条结果（显著性 / 截面 / 限制）、\\
   \footnotesize 图表家底、HEPData 是否可用};

\node[phase, right=6mm of p1] (p2)
  {\textbf{2. 为什么做}\\[2pt]
   \footnotesize 理论动机（SM 检验 / BSM 搜寻）、\\
   \footnotesize 已有约束、相对前人的进步点、\\
   \footnotesize 是在测什么物理、回答谁的预言};

\node[phase, right=6mm of p2] (p3)
  {\textbf{3. 怎么做}\\[2pt]
   \footnotesize 触发与选择、物理对象定义、\\
   \footnotesize 信号区分类、主要背景来源、\\
   \footnotesize 背景估计方法、控制区策略};

\node[phase, below=5mm of p1] (p4)
  {\textbf{4. 误差多大}\\[2pt]
   \footnotesize 主导系统误差、实验 vs 理论、\\
   \footnotesize 关联性处理、对最终结果的冲击};

\node[phase, right=6mm of p4] (p5)
  {\textbf{5. 统计与结论}\\[2pt]
   \footnotesize 拟合框架与观测量、类别划分、\\
   \footnotesize 中心值 / 置信区间 / 上限、\\
   \footnotesize SM 一致性判断\\
   \footnotesize \textcolor{cInsight!80!black}{\itshape $\to$ 并触发图表覆盖检查}};

\node[phase, right=6mm of p5] (p6)
  {\textbf{6. 批判评估}\\[2pt]
   \footnotesize 优点 / 弱点、方法学启示、\\
   \footnotesize 未来方向、可复现性评估、\\
   \footnotesize 给下一篇实验论文的借鉴};

\node[fit=(p1)(p2)(p3)(p4)(p5)(p6), inner sep=6pt, draw=cPhase!40,
      line width=0.4pt, dashed, rounded corners=6pt] (phasegroup) {};

% ---------- 右：产出 ----------
\node[product, right=16mm of p3, yshift=-8mm, text width=58mm] (output)
  {\textbf{产出：一份可机读的论文分析报告}\\[2pt]
   \footnotesize 每个问题 $\to$ 结构化 JSON\\
   \footnotesize \quad（每字段都是非空、可机读、带出处）\\
   \footnotesize $+$ \textbf{图表覆盖清单}：哪些物理图被讨论、\\
   \footnotesize \quad 哪些遗漏，强制补一轮\\
   \footnotesize $+$ \textbf{可审计工具调用轨迹}：读了哪段、\\
   \footnotesize \quad 查了哪张 HEPData 表\\
   \footnotesize $+$ 人读的 ANALYSIS.md 摘要};

% ---------- 输入 → 六维分析 → 产出 ----------
\draw[fatarrow] (paper.east) -- ++(8mm,0) |- ([xshift=-3mm]phasegroup.west);
\draw[fatarrow] ([xshift=3mm]phasegroup.east) -| ++(8mm,0) -- (output.west);

% ---------- 底部：HEP 意义 + AI 意义 ----------
\node[banner, below=12mm of phasegroup, xshift=-62mm, text width=118mm] (hepmean) {%
  {\large\textbf{\textcolor{cInsight!80!black}{HEP 物理意义}}}\\[3pt]
  \small\textbf{把"读一篇论文"这件事变成可比较的结构化数据。}\\[2pt]
  \footnotesize\ $\bullet$ 一篇实验论文本来需要博士生花 2--3 天通读 $+$ 做笔记 $+$ 核对 HEPData，\\
  \footnotesize\phantom{$\bullet$}\ 才能回答"它主导系统误差是什么""和前人工作比有什么提升"等问题\\
  \footnotesize\ $\bullet$ 现在得到一份六维结构化 JSON：每个关键数值都带出处（§/Fig/Table 锚点）\\
  \footnotesize\ $\bullet$ 下游应用：跨论文的系统误差对齐、限制对齐、方法学比较；\\
  \footnotesize\phantom{$\bullet$}\ 为 HEP-ML / 自动化 literature review / 未来对撞机唯象提供结构化前置\\
  \footnotesize\ $\bullet$ 给学生提供规范化的分析模板：不会漏系统误差、不会只抄摘要
};

\node[banner, right=10mm of hepmean, text width=118mm] (aimean) {%
  {\large\textbf{\textcolor{cInsight!80!black}{AI 方法论意义}}}\\[3pt]
  \small\textbf{一个领域原生的推理能力评估 $+$ 对齐基准。}\\[2pt]
  \footnotesize\ $\bullet$ 常见 LLM benchmark 靠"GPT-4 当裁判"——主观、贵、不可复现\\
  \footnotesize\ $\bullet$ 这里用 HEP 论文天然的结构约束当裁判信号：\\
  \footnotesize\phantom{$\bullet$}\ \textbf{数值必须对上}（带容差）、\textbf{引用锚点必须对上}、\textbf{必填字段必须填满}\\
  \footnotesize\ $\bullet$ 结果是\textbf{可审计、可复现、无 LLM-as-judge 漂移}的对齐信号\\
  \footnotesize\ $\bullet$ 三路并行的 teacher / student\_native / student\_aligned（详见第 2 页），\\
  \footnotesize\phantom{$\bullet$}\ 同时度量模型的\textbf{裸能力}与\textbf{带反馈时的可教性}
};

\end{tikzpicture}
\end{adjustbox}

\newpage

% ===================================================================
% PAGE 2 — 怎么测推理能力差距 & 怎么对齐弱模型
% ===================================================================
\begin{center}
{\LARGE\bfseries 怎么测"推理能力差距" \& 怎么把弱模型对齐到强模型}\\[2pt]
{\small 关键点：反馈信号\textbf{不是}来自第三个 LLM，而是来自论文结构的\textbf{规则化 diff}}
\end{center}

\vspace{2mm}

\begin{adjustbox}{max totalsize={\textwidth}{0.86\textheight}, center}
\begin{tikzpicture}[node distance=8mm and 12mm]

% ---------- 顶部：同一篇论文、同一套六维问题，跑三路 ----------
\node[paper, text width=72mm] (same)
  {\textbf{同一篇 HEP 论文 $+$ 同一套六维问题}\\[2pt]
   \footnotesize 控制变量后，只换"谁在读、带不带反馈"};

\node[teacher, below left=8mm and 4mm of same, text width=58mm] (teacher)
  {\textbf{Teacher}\\[2pt]
   \footnotesize 强推理模型（GPT-5.4, o-series）\\
   \footnotesize 独立走完六维分析\\
   \footnotesize \textbf{作用}：给出参考答案（不当裁判）};

\node[native, right=6mm of teacher, text width=58mm] (native)
  {\textbf{Student — Native}\\[2pt]
   \footnotesize 弱推理模型（DeepSeek V3.2）\\
   \footnotesize \textbf{完全不看 teacher 结果}独立作答\\
   \footnotesize \textbf{作用}：度量模型的\textbf{裸能力基线}};

\node[aligned, right=6mm of native, text width=58mm] (aligned)
  {\textbf{Student — Aligned}\\[2pt]
   \footnotesize 同一弱模型，\textbf{带反馈循环}\\
   \footnotesize 每个维度最多 3 轮，收敛即停\\
   \footnotesize \textbf{作用}：度量模型的\textbf{可教性 / 带反馈上限}};

\draw[arrow] (same.south) -- ++(0,-3mm) -| (teacher.north);
\draw[arrow] (same.south) -- ++(0,-3mm) -| (native.north);
\draw[arrow] (same.south) -- ++(0,-3mm) -| (aligned.north);

% ---------- 核心：反馈信号的真实来源 ----------
\node[signal, below=12mm of native, text width=186mm, align=left] (judge) {%
  \begin{minipage}{0.96\linewidth}
  \centering
  {\large\textbf{\textcolor{cSignal!85!black}{反馈从哪来？—— 不是另一个 LLM，是规则化 diff}}}\\[4pt]
  \end{minipage}

  \small
  对 Student 每一轮的回答，跟 Teacher 的参考回答做三项\textbf{确定性比对}（无 LLM 参与）：

  \medskip
  {\setlength{\leftmargini}{14pt}
  \begin{itemize}\setlength{\itemsep}{2pt}
   \item[\textcolor{cSignal!80!black}{$\blacktriangleright$}]\ \textbf{关键数值是否对上}：从 teacher/student 的回答里提取所有数字（带单位），
         按 5\% 相对容差做集合匹配 $\to$ 命中率 $\ge 80\%$ 算过
   \item[\textcolor{cSignal!80!black}{$\blacktriangleright$}]\ \textbf{引用锚点是否对上}：提取所有 §/Figure/Table/HEPData/arXiv 引用标签
         做集合交集 $\to$ 命中率 $\ge 70\%$ 算过
   \item[\textcolor{cSignal!80!black}{$\blacktriangleright$}]\ \textbf{必填字段是否填满}：每个维度预先声明了必须非空的字段（例如"第 4 维必须有
         \textit{dominant\_systematics}"） $\to$ 填满 $\ge 80\%$ 算过
  \end{itemize}
  }

  \medskip
  没过 $\to$ 把\textbf{差异摘要}填入中文反馈模板（只说"缺了哪些数字、哪些锚点、哪些字段"，
  \textbf{不泄露 teacher 的答案原文}，防止学生抄而不验），回塞 student 再跑一轮。
};

% ---------- Aligned 的内循环简图 ----------
\node[aligned, below left=10mm and 22mm of judge, text width=54mm] (round0)
  {\textbf{Round 0}\\[2pt]
   \footnotesize 学生裸跑（等同 Native）};

\node[signal, right=10mm of round0, text width=48mm] (diff)
  {\textbf{规则化 diff}\\[2pt]
   \footnotesize 数值 / 引用 / 必填 三项};

\node[aligned, right=10mm of diff, text width=54mm] (retry)
  {\textbf{Round 1--3}\\[2pt]
   \footnotesize 带\textbf{差异摘要}重跑\\
   \footnotesize （摘要中无 teacher 原文）};

\node[product, right=10mm of retry, text width=52mm] (best)
  {\textbf{Best-of-rounds}\\[2pt]
   \footnotesize 按打分取每维最优\\
   \footnotesize 防止 flaky 轮覆盖好答案};

\draw[arrow] (round0) -- (diff);
\draw[signalarrow] (diff) -- node[above, font=\footnotesize]{没过} (retry);
\draw[loop] (retry.south) .. controls +(0,-7mm) and +(0,-7mm) ..
  node[below, font=\footnotesize\itshape]{再来一轮（$\le$ 3）} (diff.south);
\draw[arrow] (retry) -- (best);

% ---------- 底部：最终产出 + AI 价值点 ----------
\node[product, below=10mm of round0, text width=96mm, align=left] (metrics)
  {\textbf{最终度量指标（每篇论文、每个维度）}\\[3pt]
   \footnotesize\ $\bullet$ \textbf{Native 通过率} —— 模型裸能力\\
   \footnotesize\ $\bullet$ \textbf{Aligned 收敛率} —— 带反馈后能否达标\\
   \footnotesize\ $\bullet$ \textbf{收敛所需轮数} —— 可教性的精细粒度\\
   \footnotesize\ $\bullet$ \textbf{数值 / 引用命中率} —— 对物理细节的把握度\\
   \footnotesize\ $\bullet$ \textbf{工具调用轨迹编辑距离} —— 读论文的"路径"像不像老手};

\node[insight, right=10mm of metrics, text width=108mm, align=left] (value)
  {\textbf{\textcolor{cInsight!80!black}{这个设计的 AI 方法论价值}}\\[3pt]
   \footnotesize\ $\bullet$ \textbf{反馈信号可复现 $+$ 便宜} —— 不用 GPT-4 当裁判，不受其主观性 / API 漂移影响\\
   \footnotesize\ $\bullet$ \textbf{领域结构先验作信号} —— "HEP 论文必须报什么"本身是强约束，信号比通用打分更锐利\\
   \footnotesize\ $\bullet$ \textbf{不泄露答案的反馈} —— 学生被迫\textbf{自己去论文里验证}，而非抄 teacher；\\
   \footnotesize\phantom{$\bullet$}\ 这使 aligned 过程测的是"带指引的推理能力"而非"抄写能力"\\
   \footnotesize\ $\bullet$ \textbf{同时度量裸能力与可教性} —— 给弱推理模型的"潜力上限"一个定量刻度\\
   \footnotesize\ $\bullet$ 方法本身\textbf{领域可迁移}：只要领域能给出"必答问题 $+$ 可抽取的数值 / 引用"，就能套用
};

\draw[darrow] (best.south) .. controls +(0,-10mm) and +(0,10mm) .. (value.north);
\draw[darrow] (round0.south) -- (metrics.north);

\end{tikzpicture}
\end{adjustbox}

\end{document}
